\section{Results}

\subsection{Tongji Multi-Seed Results}

Table~\ref{tab:tongji_multiseed} reports bidirectional Tongji averages over seeds 42, 2026, and 2705. B6 outperforms both the B1 CE + SupCon baseline and the B2 fixed-Gabor variant on the main identification and verification metrics.

\begin{table*}[t]
\centering
\caption{Tongji bidirectional multi-seed results. Values are mean +/- standard deviation over seeds 42, 2026, and 2705.}
\label{tab:tongji_multiseed}
\resizebox{\textwidth}{!}{%
\begin{tabular}{llccccccccc}
\toprule
Method & Protocol & Rank-1 & Rank-5 & Macro-F1 & EER & TAR@1e-2 & TAR@1e-3 & Time & Params & FLOPs \\
\midrule
B1 ResNet18 + CE + SupCon & Bidir. Avg. & 94.37 +/- 0.88 & 96.24 +/- 0.67 & 93.61 +/- 1.05 & 1.95 +/- 0.30 & 96.92 +/- 0.66 & 91.56 +/- 1.91 & 1.75 +/- 0.07 ms & 11.46M & 1.82G \\
B2 FixedGaborResNet18 & Bidir. Avg. & 93.71 +/- 0.55 & 95.60 +/- 0.26 & 92.81 +/- 0.60 & 2.09 +/- 0.21 & 96.90 +/- 0.37 & 91.78 +/- 0.74 & 4.87 +/- 1.71 ms & 11.82M & 2.71G \\
B6 ResNet18 + BNNeck + ArcFace & Bidir. Avg. & 96.07 +/- 0.73 & 97.45 +/- 0.59 & 95.56 +/- 0.84 & 1.45 +/- 0.30 & 98.23 +/- 0.53 & 95.20 +/- 1.21 & 4.89 +/- 0.81 ms & 11.46M & 1.82G \\
\bottomrule
\end{tabular}%
}
\end{table*}

Relative to B1, B6 improves Rank-1 by +1.70 percentage points, reduces EER by -0.51 percentage points, and improves TAR@FAR=1e-3 by +3.63 percentage points. These gains support the scoped cross-session claim for BNNeck + ArcFace.

Relative to B2, B6 improves Rank-1 by +2.36 percentage points and TAR@FAR=1e-3 by +3.41 percentage points. B6 also uses 0.36M fewer parameters and 0.89G fewer FLOPs than B2 in the recorded summary. This makes B6 the stronger final method direction than the earlier fixed-Gabor variant.

\subsection{Direct Comparisons}

\begin{table}[t]
\centering
\caption{Bidirectional 3-seed average deltas. Positive values are better for Rank-1, Rank-5, Macro-F1, and TAR. Negative values are better for EER, time, parameters, and FLOPs.}
\label{tab:direct_comparison}
\resizebox{\columnwidth}{!}{%
\begin{tabular}{lccccccccc}
\toprule
Comparison & Rank-1 & Rank-5 & Macro-F1 & EER & TAR@1e-2 & TAR@1e-3 & Time & Params & FLOPs \\
\midrule
B6 - B1 & +1.70 & +1.21 & +1.95 & -0.51 & +1.31 & +3.63 & +3.14 ms & +0.00M & +0.00G \\
B6 - B2 & +2.36 & +1.85 & +2.75 & -0.65 & +1.33 & +3.41 & +0.03 ms & -0.36M & -0.89G \\
\bottomrule
\end{tabular}%
}
\end{table}

\subsection{IITD Secondary Validation}

Table~\ref{tab:iitd} reports the IITD within split. B6 remains competitive and improves over B2, but B1 remains the strongest method on this split. This result is important for scope control: the paper should not claim universal superiority.

\begin{table}[t]
\centering
\caption{IITD within split seed42 results.}
\label{tab:iitd}
\resizebox{\columnwidth}{!}{%
\begin{tabular}{lcccccc}
\toprule
Method & Rank-1 & Rank-5 & Macro-F1 & EER & TAR@1e-2 & TAR@1e-3 \\
\midrule
B1 ResNet18 + CE + SupCon & 99.13 & 99.57 & 98.84 & 0.36 & 99.70 & 99.41 \\
B2 FixedGaborResNet18 & 98.26 & 99.13 & 97.68 & 0.71 & 99.29 & 98.93 \\
B6 ResNet18 + BNNeck + ArcFace & 98.48 & 99.13 & 97.97 & 0.66 & 99.41 & 99.11 \\
\bottomrule
\end{tabular}%
}
\end{table}
